\section{Exactly Solving for the Ground State}
According to equation (29) of Ref.~\cite{Kitaev_2006}, we can diagonalize the Kitaev Hamiltonian in the momentum space if the gauge is fixed to be $u_{ij} = +1\forall\; (ij)$, 
\[\hat{H} = \sum_{\bo{k}}\Psi^\dagger_{k}U(k)\Psi_{k} = \frac{1}{2}\sum_{\bo{k}}\lrm{is_{\bo{k}}d_{k,A}^\dagger d_{k,B} - is_{\bo{k}}^*d_{k,B}^\dagger d_{k,A}}\]
where
\begin{equation}\label{eq:spinor-def}
    \Psi_k \equiv \frac{1}{\sqrt{N/2}}\sum_{\bo{r}}e^{-ik\cdot \bo{r}}\begin{pmatrix}
        c_{\bo{r},A} \\
        c_{\bo{r},B}
       \end{pmatrix} 
       = 
       \begin{pmatrix}
         d_{k,A} \\
         d_{k,B}
       \end{pmatrix}
\end{equation}
\[U(k) = \frac{1}{2}\begin{pmatrix}
                      0 & is_{\bo{k}} \\
                      -is_{\bo{k}}^* & 0 
                    \end{pmatrix} \quad , \quad s_{\bo{k}} = 2\lr{J_xe^{ik\cdot \lr{\bo{a}_2 - \bo{a}_1}} + J_ye^{ik\cdot \bo{a}_2} + J_z}\]
Note that the index $\bo{r}$ above is the Bravais lattice point, and $A$,$B$ represent the sublattice. If we further define a normal fermion operator, 
\begin{equation}\label{gamm def.e}
  \gamma_{\bo{r}} \equiv \frac{1}{2}\lr{c_{{\bo{r}},A} + ic_{{\bo{r}},B}}
\end{equation}
 then its Fourier transform would be 
\[\gamma_{k} = \frac{1}{\sqrt{N/2}}\sum_{{\bo{r}}}\gamma_{{\bo{r}}}e^{-ik\cdot \bo{r}} = \frac{1}{2}\lr{d_{k,A} + id_{k,B}}\]
In this way, the operator $d$ can be expressed in terms of $\gamma$ fermion,
\begin{equation}\label{eq:d-gamma-transform}
    \Psi_{\bo{k}}
    =
    \begin{pmatrix}
        d_{\bo{k}A} \\
        d_{\bo{k}B}
    \end{pmatrix}
    =
    \begin{pmatrix}
        1 & 1 \\
        -i & i
    \end{pmatrix}
    \begin{pmatrix}
        \gamma_{\bo{k}} \\
        \gamma^\dagger_{-\bo{k}}
    \end{pmatrix}
\end{equation}
Using the $\gamma$ operators to rewrite the Hamiltonian gives
\[\hat{H} = \sum_{\bo{k}}\lrm{R_\bo{k}\cdot\lr{\gamma_{\bo{k}}^\dagger \gamma_{\bo{k}} - \gamma_{-\bo{k}}\gamma_{-\bo{k}}^\dagger} - iI_\bo{k}\cdot\lr{\gamma_{k}^\dagger \gamma_{-\bo{k}}^\dagger - \gamma_{-\bo{k}}\gamma_{k}}} = \sum_{\bo{k}}\phi_k^\dagger\bo{H}(k)\phi_k\]
where $R_{\bo{k}} \equiv \mrm{Re}\lrg{s_{\bo{k}}}$ and $I_k \equiv \mrm{Im}\lrg{s_{\bo{k}}}$, and 
\[\bo{H}(k) \equiv 
\begin{pmatrix}
    R_{\bo{k}} & -iI_{\bo{k}} \\
    iI_k & -R_{\bo{k}} 
\end{pmatrix} 
\quad,\quad 
\phi_k 
\equiv 
\begin{pmatrix}
    \gamma_{\bo{k}} \\
    \gamma_{-\bo{k}}^\dagger 
\end{pmatrix}
=
\frac{1}{2}
\begin{pmatrix}
    1 & i \\
    1 & -i
\end{pmatrix}
\Psi_{\bo{k}}
\]
Since 
\[\bo{H}(k)^2 = \lr{R_{\bo{k}}^2 + I_k^2}\mathds{1} = \abs{s_{\bo{k}}}^2\mathds{1}\]
we can see that the eigenvalues are $\epsilon_{k} = \pm\abs{s_{\bo{k}}}$, which is consistent with Kitaev's result (equation (32) in Ref.\cite{Kitaev_2006}). To diagonalize the matrix $\bo{H}$ we have to apply the Bogoliubov transformation,
\begin{align*}
U &= \frac{1}{\sqrt{2\abs{s_{\bo{k}}}}}
        \begin{pmatrix}
           \mrm{sgn}(I_k)\sqrt{\abs{s_{\bo{k}}}+R_{\bo{k}}} & -i\sqrt{\abs{s_{\bo{k}}}-R_{\bo{k}}} \\
           i\sqrt{\abs{s_{\bo{k}}}-R_{\bo{k}}} & -\mrm{sgn}(I_k)\sqrt{\abs{s_{\bo{k}}}+R_{\bo{k}}} 
         \end{pmatrix} 
         = 
         \begin{pmatrix}
           \mrm{sgn}(I_k) u_k & -iv_k \\
           i v_k &  -\mrm{sgn}(I_k) u_k
         \end{pmatrix} \\
         &= 
         v_k\sigma^y + \mrm{sgn}(I_k)u_k\sigma^z
\end{align*}
such that\footnote{Those parameters in the matrix $U$ have the following properties: \textbf{(i)} $u_k = u_{-k}$, \textbf{(ii)} $v_{k} = v_{-k}$, and  \textbf{(iii)} $\mrm{sgn}(I_{k}) = -\mrm{sgn}(I_{-k})$} 
\[U\bo{H}U^\dagger = \begin{pmatrix}
                       \abs{s_{\bo{k}}} & 0 \\
                       0 & -\abs{s_{\bo{k}}} 
                     \end{pmatrix}\]
Therefore, the excitation that diagonalize the Hamiltonian is 
\begin{equation}\label{Bogo.e}
\begin{pmatrix}
    a_k \\
    a_{-k}^\dagger 
  \end{pmatrix} 
  = 
  U
    \begin{pmatrix}
     \gamma_{\bo{k}} \\
     \gamma^\dagger_{-k} 
    \end{pmatrix} 
    = 
    \begin{pmatrix}
        \mrm{sgn}(I_k) u_k\gamma_{\bo{k}} - iv_k\gamma^\dagger_{-k} \\
         i v_k\gamma_{\bo{k}} - \mrm{sgn}(I_k) u_k\gamma^\dagger_{-k} 
   \end{pmatrix}
\end{equation}
As a result, the Kitaev ground state is the vacuum state of the Bogoliubov fermion, i.e., $a_{k}\ket{M} = 0$, which means $\ket{M}$ can be written as a BCS ground state just as we met in the BCS theory. The ground state can be written as the coherent state\footnote{
The coefficient $g_k$ can be represented in terms of $u_k$ and $v_k$. As I have mentioned, the BCS-like ground state is expected to be the vacuum state of $a_k$ fermions for all possible $k$, thus, we have
\[
\lr{\mrm{sgn}(I_k) u_k\gamma_{\bo{k}} - iv_k\gamma_{-\bo{k}}^\dagger}\prod_{k}\lr{\mathds{1} + g_k\gamma_{\bo{k}}^\dagger \gamma_{-\bo{k}}^\dagger}\ket{\vac_\gamma} = 0
\]
Based on this, we obtain that 
\[\mrm{sgn}(I_k) u_kg_k - iv_k = 0\]
therefore
\begin{equation}\label{g_k solution.e}
  g_k = \mrm{sgn}(I_k)i\frac{v_k}{u_k} = \mrm{sgn}(I_k) i\sqrt{\frac{\abs{s_{\bo{k}}} - R_{\bo{k}}}{\abs{s_{\bo{k}}} + R_{\bo{k}}}} = i\frac{\abs{s_{\bo{k}}} - R_{\bo{k}}}{I_k}
\end{equation}
} of the $\gamma$-fermion pair,
\begin{align*}
  \ket{M} &= A\lrm{g}\exp\lr{\sum_{\bo{k}}g_k \gamma^\dagger_k \gamma^\dagger_{-k}}\ket{\vac_\gamma}\\
   &= A\lrm{g}\exp\lr{\sum_{{\bo{r}}{\bo{r}'}}g_{{\bo{r}}{\bo{r}'}} \gamma^\dagger_{\bo{r}} \gamma^\dagger_{{\bo{r}'}}}\ket{\vac_\gamma}\\
   &= A\lrm{g}\lrm{\mathds{1} + \sum_{{\bo{r}}{\bo{r}'}}g_{{\bo{r}}{\bo{r}'}} \gamma^\dagger_{\bo{r}} \gamma^\dagger_{{\bo{r}'}} + \frac{1}{2!}\sum_{{\bo{r}}{\bo{r}'}\kappa\lambda}g_{{\bo{r}}{\bo{r}'}}g_{\kappa\lambda
   } \gamma^\dagger_{\bo{r}} \gamma^\dagger_{{\bo{r}'}}\gamma^\dagger_\kappa \gamma^\dagger_{\lambda} + \cdots}\ket{\vac_\gamma}
\end{align*}
where $\ket{\vac_\gamma}$ is the vacuum state of the $\gamma$-fermions, and $g_{{\bo{r}}{\bo{r}'}} = \frac{1}{2N}\sum_{\bo{k}}g_ke^{-ik\cdot(\bo{r}-\bo{r}')}$. 

\subsection{With External Magnetic Field}
For simplicity, we focus only on the isotropic limit~\cite{Kitaev_2006}. In this situation, applying an external magnetic field $\bo{h} = \lr{h_x,h_y,h_z}$, which is equivalent to adding a Zeeman perturbation $H' = -\sum_{i}\bo{h}\cdot\bog{\sigma}$, makes the Kitaev model equivalent to the Haldane honeycomb model~\cite{Haldane_1988} up to the third-order perturbation. The Hamiltonian under this approximation in the Majorana fermion representation reads
\begin{align*}
    H 
    &= 
    \frac{1}{2}J\sum_{\bo{r}}
    \lr{
        ic_{\bo{r}A}c_{\bo{r}+\hat{\bo{n}}_1,B} 
        + 
        ic_{\bo{r}A}c_{\bo{r}+\hat{\bo{n}}_2,B}
        +
        ic_{\bo{r}A}c_{\bo{r}B}
        +\mrm{H.c.}
    } \\
    &\quad\quad\quad
    + 
    \frac{1}{2}\kappa
    \sum_{\bo{r}}\sum_{\mu = A,B}
    (-1)^\mu
    \lr{-i
        c_{\bo{r}\mu}c_{\bo{r}+\hat{\bo{n}}_1,\mu}
        +i
        c_{\bo{r}\mu}c_{\bo{r}+\hat{\bo{n}}_2,\mu}
        -i
        c_{\bo{r}\mu}c_{\bo{r}+\hat{\bo{n}}_1-\hat{\bo{n}}_2,\mu}
        +
        \mrm{H.c.}
    }
\end{align*}
Expand the Majorana fermions in terms of the momentum modes $c_{\bo{r}\mu} = \frac{1}{\sqrt{N/2}}\sum_{\bo{k}} e^{i\bo{k}\cdot\bo{r}} d_{\bo{k}\mu}$, the second term proportional to $\kappa$ becomes
\[
\kappa
\sum_{\bo{k}}
\sum_{\mu}
(-1)^{\mu}
\frac{1}{2}
\lr{
    -i e^{i\bo{k}\cdot\hat{\bo{n}}_1}
    +i e^{i\bo{k}\cdot\hat{\bo{n}}_2}
    -i e^{i\bo{k}\cdot\lr{\hat{\bo{n}}_1 - \hat{\bo{n}}_2}}
    + \mrm{c.c.}
}   
d^\dagger_{\bo{k}\mu} d_{\bo{k}\mu}
\equiv
\frac{1}{2}
\sum_{\bo{k}}
\Delta_{\bo{k}}
\lr{
d^\dagger_{\bo{k}A} d_{\bo{k}A}
-
d^\dagger_{\bo{k}B} d_{\bo{k}B}
}
\]
where
\begin{equation}\label{eq:gap-def}
    \Delta_{\bo{k}} \equiv
    2\kappa
    \lrm{
    \sin\lr{\bo{k}\cdot\hat{\bo{n}}_1}
    -
    \sin\lr{\bo{k}\cdot\hat{\bo{n}}_2}
    +
    \sin\lr{\bo{k}\cdot\lr{\hat{\bo{n}}_1 - \hat{\bo{n}}_2}}
    }
\end{equation}
Converting the Hamiltonian into a spinor form defined by Eq.~\eqref{eq:spinor-def}, we obtain
\begin{align*}
H 
&= 
\frac{1}{2}
\sum_{\bo{k}}
\Psi_{\bo{k}}^\dagger
\begin{pmatrix}
  \Delta_{\bo{k}} & is_{\bo{k}} \\
  -is_{\bo{k}}^* & -\Delta_{\bo{k}}
\end{pmatrix}
\Psi_{\bo{k}} \\
&=
\frac{1}{2}
\sum_{\bo{k}}
\phi_{\bo{k}}^\dagger
\begin{pmatrix}
    1 & i \\
    1 & -i
\end{pmatrix}
\begin{pmatrix}
  \Delta_{\bo{k}} & is_{\bo{k}} \\
  -is_{\bo{k}}^* & -\Delta_{\bo{k}}
\end{pmatrix}
\begin{pmatrix}
    1 & 1 \\
    -i & i
\end{pmatrix}
\phi_{\bo{k}} \\
&= 
\sum_{\bo{k}}
\phi_{\bo{k}}^\dagger
\begin{pmatrix}
    R_{\bo{k}} & \Delta_{\bo{k}} - iI_{\bo{k}} \\
    \Delta_{\bo{k}} + iI_{\bo{k}} & -R_{\bo{k}}
\end{pmatrix}
\phi_{\bo{k}}
\equiv 
\sum_{\bo{k}} 
\phi_{\bo{k}}^\dagger
\bo{H}(\bo{k})
\phi_{\bo{k}}
\end{align*}
The BdG Hamiltonian represented in terms of Pauli matrices reads
\begin{equation}\label{eq:BdG-magnetic}
\bo{H}(\bo{k})
=
\Delta_{\bo{k}} \tau^x
+
I_{\bo{k}} \tau^y
+
R_{\bo{k}} \tau^z
\end{equation}
\[
\epsilon(\bo{k})
=
\sqrt{
\abs{s_{\bo{k}}}^2 + \Delta_{\bo{k}}^2
}
\]
The Bogoliubov transformation matrix is
\begin{align}
    U 
    &= 
    \frac{1}{\sqrt{2 \epsilon(\bo{k}) \lr{\epsilon(\bo{k}) + R_{\bo{k}}}}}
    \begin{pmatrix}
        \epsilon(\bo{k}) + R_{\bo{k}} & \Delta_{\bo{k}} - iI_{\bo{k}} \\
        \Delta_{\bo{k}} + iI_{\bo{k}} & -\epsilon(\bo{k}) - R_{\bo{k}}
    \end{pmatrix} \nonumber\\
    & =
    \frac{1}{\sqrt{2\epsilon(\bo{k})}}
    \begin{pmatrix}
        \sqrt{\epsilon(\bo{k}) + R_{\bo{k}}} & e^{-i\varphi_{\bo{k}}} \sqrt{\epsilon(\bo{k}) - R_{\bo{k}}} \\
        e^{i\varphi_{\bo{k}}} \sqrt{\epsilon(\bo{k}) - R_{\bo{k}}} & - \sqrt{\epsilon(\bo{k}) + R_{\bo{k}}} 
    \end{pmatrix} \nonumber \\
    &\equiv
    \begin{pmatrix}
        u_{\bo{k}} & e^{-i\varphi_{\bo{k}}} v_{\bo{k}} \\
        e^{i\varphi_{\bo{k}}}v_{\bo{k}} & -u_{\bo{k}}
    \end{pmatrix} = U^\dagger = U^{-1}
\end{align}
where $\varphi_{\bo{k}} \equiv \arg\lr{\Delta_{\bo{k}} + iI_{\bo{k}}}$. Therefore, the elementary excitations that diagonalize the Hamiltonian are
\begin{equation}\label{eq:magnetic-field-elementary-excitation}
\begin{pmatrix}
    a_k \\
    a_{-k}^\dagger 
  \end{pmatrix} 
  = 
  U
    \begin{pmatrix}
     \gamma_{\bo{k}} \\
     \gamma^\dagger_{-k} 
    \end{pmatrix} 
    = 
    \begin{pmatrix}
        u_{\bo{k}}\gamma_{\bo{k}} + e^{-i\varphi_{\bo{k}}} v_k \gamma_{-\bo{k}}^\dagger \\
        e^{i\varphi_{\bo{k}}} v_k \gamma_{\bo{k}} + u_{\bo{k}} \gamma_{-\bo{k}}^\dagger
   \end{pmatrix}
\end{equation}

